% !PRINTED 235
\weberpage{235}
%<<CONT
and by (4) this equation can also be represented thus:
\weq{7}{m\,(a\,u + b\,v) = (a'\,u + b'\,v)^2 - c\,(a\,u + b\,v)^2.}
Now set
\weq{8}{y = \frac{a'\,u + b'\,v}{a\,u + b\,v},}
and form the equation of the fifth degree whose roots are the five values of $y$:
\weq{9}{y^5 + c_1\,y^4 + c_2\,y^3 + c_3\,y^2 + c_4\,y + c_5 = 0;}
in it the coefficients $c$ can be determined more closely in the following way.

From (8) we derive
\[
y + \sqrt{c} = \frac{a'\,u + b'\,v + \sqrt{c}\,(a\,u + b\,v)}{a\,u + b\,v}
\]
and from this, by (7),
\[
y + \sqrt{c} = \frac{m}{a'\,u + b'\,v - \sqrt{c}\,(a\,u + b\,v)},
\]
\[
\frac{m}{y + \sqrt{c}} = a'\,u + b'\,v - \sqrt{c}\,(a\,u + b\,v).
\]
But from this it follows, by (3), that
\[
S\left[\frac{1}{y + \sqrt{c}}\right] = 0, \quad S\left[\frac{1}{(y + \sqrt{c})^2}\right] = 0.
\]
If, then, we derive from (9) the equation whose roots are the values
\[
\xi = \frac{1}{y + \sqrt{c}}
\]
that is, if we set
\[
y = \frac{1 - \xi\,\sqrt{c}}{\xi}
\]
an equation for $\xi$ must arise in which the coefficients of $\xi^4$ and $\xi^3$ vanish, and indeed whichever sign we give to the square root $\sqrt{c}$. In this way one obtains four equations between the coefficients $c_\nu$ of equation (9).

The equation for $\xi$ becomes, namely,
\[
(1 - \xi\,\sqrt{c})^5 + c_1\,\xi\,(1 - \xi\,\sqrt{c})^4 + c_2\,\xi^2\,(1 - \xi\,\sqrt{c})^3 + c_3\,\xi^3\,(1 - \xi\,\sqrt{c})^2 + c_4\,\xi^4\,(1 - \xi\,\sqrt{c}) + c_5\,\xi^5 = 0.
\]
If one here sets the coefficients of $\xi^4$ and $\xi^3$ equal to 0, it follows that
%CONT>>
